\chapter{現段階の結論}

\begin{table}[H]
\centering
\caption{能力別の到達状況}
\begin{tabularx}{\textwidth}{>{\bfseries}p{35mm}p{30mm}X}
\toprule
能力 & 判定 & 根拠と境界\\
\midrule
データ生産・管理 & 達成 & 51有効プロジェクト、2,413軌跡、専用収集・編集基盤\\
本番運用モデルの学習・保存 & 構築済み & 25確定版、1,051軌跡、59,990ステップ履歴、学習済みモデルを確認\\
双腕全工程の現場統合 & 現場で統合動作を確認 & 把持、開栓、本体・キャップの分別投入を現場で観測\\
長時間反復 & 観測済み & 1時間超・約2本/分は現場概算で、標準手順による計測ではない\\
多条件への初期適応 & 一部達成 & 多条件データと現場兆候あり、固定試験は未実施\\
キャップなし判断・向き判断 & 継続改善 & 空回し、逆向き把持を観測し、対応データ入口を特定\\
強化学習による実機向上 & 未確認 & オフライン閉ループと約6.7\%改善のみ\\
標準成功率・処理量 & 未測定 & 完全映像、逐次台帳、固定試験条件が必要\\
\bottomrule
\end{tabularx}
\end{table}

\section{初年度に確立したもの}

本プロジェクトは、専用収集、データ編集、確定版公開、大規模学習、学習済みモデル、現場運用、実行監査までを接続した。本番学習には1,051軌跡、879,852フレームを用い、学習履歴は59,990ステップ、約51.4時間に及ぶ。現場では双腕によるボトル把持、協調開栓、ボトル本体とキャップの分別投入という全工程を観測した。

この成果は、単発の把持モデルではなく、視覚判断、双腕協調、接触、分離、分類、復帰を含む多段階タスクの実機システムを一体化した点にある。さらに、41回の学習試行、8,223件のアテンション記録、835軌跡の強化学習研究が、次年度の比較改善に使える資産として残っている。

\section{現段階の最良結果と表現範囲}

最良の現場結果は、全タスク反復、1時間超の連続稼働、平均約2本/分という観測である。ただし完全映像と計数台帳がないため、後二者は概算として報告する。標準試験手順に基づく成功率、条件別成功率、処理量の信頼区間は未記録であり、学習損失から代用しない。

\resultbox{現場で全工程と長時間反復を示したこと、ならびにデータから運用までの再利用可能な閉ループを構築したことが初年度の中核成果である。次年度は、この実働能力を固定条件・逐次記録・工程別判定で定量化し、キャップの有無と把持方向への適応を強化する。}

\section{主な改善余地}

キャップのないボトルでも開栓動作へ進む現象と、一部を逆向きに把持する現象がある。前者については、キャップなし158軌跡の指示文に「開栓を省略」と明示した例がないことを確認した。後者については方向404軌跡と反転127軌跡が既にあり、次は口部方向と把持端の結果ラベルを加える。いずれも問題発見だけでなく、追加データと対照試験へ直結する具体的な改善経路が得られている。

\section{安定展示への判断}

システムは、技術展示として全工程を見せる段階には到達している。より厳格な「再現性のある安定展示」を宣言するには、固定初期条件、逐次計数、失敗類型、復旧時間、標準試験手順に基づく成功率を記録する必要がある。したがって次年度は、能力向上と同時に標準評価基盤を継続運用できる形に整備することが最優先となる。
